\section{Experimental Evaluation}

The proposed algorithm is studied by performing several experiments to explore its properties such as convergence, robustness, and behavior under various graph structures. Unless stated otherwise, all experiments have been performed on the graph presented in Section~2.

\subsection{Convergence Behavior}

First, we study the convergence properties of the proposed propagation algorithm to determine whether it converges to a stable reliability distribution. In order to achieve this, we initialize all sources with equal starting reliability and apply the propagation algorithm described in Section~3 iteratively until the convergence condition is met. We measure the maximum absolute difference between the source reliability vectors of two consecutive iterations during the propagation process.

We can see in Figure 2 that the maximum difference in source reliability values diminishes at each subsequent iteration up to the point where the convergence condition is met.

\begin{figure}[t]
    \centering
    \includegraphics[width=0.75\linewidth]{figures/convergence.png}
    \caption{Maximum change between consecutive source reliability vectors during propagation. The reliability vector converges after 37 iterations.}
    \label{fig:convergence}
\end{figure}

\subsection{Initialization Independence}

Having demonstrated convergence of the propagation algorithm, we now explore whether the resulting reliability distribution is independent of the initial distribution. To perform this analysis, we test two methods of initialization. The first method assigns each source equal initial reliability. The second method gives each source a random reliability value, which is then normalized.

We observe that, despite starting at different points, both methods lead to the same reliability distribution. This suggests that the converged solution may be independent of the initial distribution. The resulting reliability distribution is therefore stable and consistent.

\subsection{Graph Connectivity}

Having explored the convergence properties of the propagation algorithm, we now turn our attention to the effect that graph connectivity has on reliability propagation. Since the process of reliability propagation occurs alongside assertion edges between source nodes and their claims, graph connectivity impacts which sources can propagate their reliability to one another.

For the purpose of studying the dynamics of this phenomenon, we first consider the toy graph shown in Figure~\ref{fig:connectivity}.

\begin{figure}[t]
    \centering
    \includegraphics[width=0.82\linewidth]{figures/connectivity.png}
    \caption{Toy graph consisting of two connected components. Sources $A$ and $B$ share claims $C1$ and $C2$, while source $D$ is isolated through claims $C3$ and $C4$. Reliability propagates only within connected components.}
    \label{fig:connectivity}
\end{figure}

It turns out that reliability propagates only among the connected components of the graph. Sources disconnected from the rest of the network do not receive reliability from outside and thus converge independently of the rest of the graph.

This example clearly shows that graph connectivity plays an important part in assertion evaluation. The degree to which the sources share claims determines the pattern of reliability propagation across the network.

\subsection{Claim Representation}

In view of the above observations regarding the effect of graph connectivity on reliability propagation, let’s examine the effects that the choice of representation of claims has on the graph structure. As the graph is constructed from the set of claims and assertions, the way claims are represented affects the graph structure in a significant way.

In particular, initially, claims were considered as unique combinations made up of an entity, its attribute, and the observed value of the attribute. While such a representation kept track of all observed values, it led to a highly disconnected graph in which many claims were asserted by only one source. Therefore, the reliability propagation was restricted to small connected components of the graph.

To address this problem, we redefine a claim as an ordered pair of the entity and attribute alone. Using this definition, multiple sources asserting different values for the same attribute get connected through a common claim node instead of being kept separate as distinct nodes.

We find this way of representing claims causes a significant increase in the connectivity of the graph. Reliability can now spread through a larger part of the graph. Consequently, conflicting assertions form a single evaluation path, thus allowing us to estimate reliability from the graph structure as a whole.

\subsection{Agreement-Weighted Propagation}

After improving the connection of our graph, we explore whether the propagation algorithm can exploit the level of agreement between contradicting assertions. In the original propagation algorithm, all the assertions carried equal influence during the reliability propagation regardless of the popularity of the asserted value.

\begin{figure}[H]
    \centering
    \includegraphics[width=0.78\linewidth]{figures/agreement_weights.png}
    \caption{Agreement-weighted propagation. Four sources assert competing values for the same claim. Assertions supporting the most frequently asserted value receive greater propagation weight than assertions supporting less common values.}
    \label{fig:agreement}
\end{figure}

This is where agreement-weighting comes into play. The assertions on each claim are weighted depending on the percentage of the sources asserting that particular value. Assertions which support more popular values are given higher weights compared to those which assert rare values, as pictured in Figure~\ref{fig:agreement}.

We have seen that using agreement weighting enables us to incorporate local agreement information into the propagation process while retaining the convergence properties of the algorithm. As opposed to the original process where all the assertions were treated equally, we are now considering the distribution of the competing values on the same claim.
